\exercisegroup

\begin{enumerate}
\def\labelenumi{\arabic{enumi}.}
\item
  设 E 和 F 为两个集合，设 f 为 E 到 F 的一个单射，并设 g 为 F 到 E 的一个映射。证明存在 E 的两个子集 A、B，使得 B = E - A，以及 F 的两个子集 A'、B'，使得 B' = F - A'，并且满足 A' = f(A) 且 B = g(B')。（设 R = E - g(F)，并令 h = g∘f；取 A 为 E 的满足 M ⊃ R ∪ h(M) 的那些子集 M 的交集。）
\item
  如果E和F是不同的集合，证明 \(E^{F} \neq F^{E}\) 。推导出，如果 E 和 F 是基数 2 和 \(4(=2+2)\) ，那么集合 \(E^{F}, F^{E}\) 中至少有一个不是基数。
\item
  让 \((\mathfrak{a}_t)_{t\in I}\) 和 \((\mathfrak{b}_t)_{t\in I}\) 是两族基数，使得 \(\mathfrak{b}_t\geqslant 2\) 对于每一个 \(t\in I\) .
\end{enumerate}

\begin{enumerate}
\def\labelenumi{(\alph{enumi})}
\tightlist
\item
  证明，如果 \(a_{\iota} \leqslant b_{\iota}\) 对于每一个 \(\iota \in I\) ，那么
\end{enumerate}

\[
\sum_ {i \in I} a _ {i} \leqslant \underset {i \in I} {\mathbf {P}} b _ {i}.
\]

\begin{enumerate}
\def\labelenumi{(\alph{enumi})}
\setcounter{enumi}{1}
\tightlist
\item
  证明：如果 \(\mathfrak{a}_{\iota} < \mathfrak{b}_{\iota}\) 对于每一个 \(\iota \in \mathbf{I}\) ，那么
\end{enumerate}

\[
\sum_ {i \in I} a _ {i} <   \underset {i \in I} {\mathbf {P}} b _ {i}.
\]

(注意，一个乘积 \(\prod_{i\in I}E_i\) 不能是一个族 \((A_i)_{i\in I}\) 的并集，使得 Card \((A_i) < \text{Card}(E_i)\) 对于所有 \(i\in I\) ，通过观察 Card \((\text{pr}_i(A_i)) < \text{Card}(E_i).\)

\begin{enumerate}
\def\labelenumi{\arabic{enumi}.}
\setcounter{enumi}{3}
\tightlist
\item
  令 E 为一个集合，并令 f 为从 \(\mathfrak{P}(\mathbf{E}) - \{\varnothing\}\) 到 E 中的一个映射，使得对于 E 的每个非空子集 X，我们有 \(f(\mathbf{X}) \in \mathbf{X}\) （“选择函数”）。
\end{enumerate}

\begin{enumerate}
\def\labelenumi{(\alph{enumi})}
\item
  设 \(\mathfrak{b}\) 为一个基数，并令 \(\mathbf{A}\) 为 \(x \in \mathbf{E}\) 中所有使得 \(\operatorname{Card}(\overline{f}^1(x)) \leqslant \mathfrak{b}\) 的元素所成的集合 . 证明：如果 \(\mathfrak{a} = \operatorname{Card}(\mathbf{A})\) ，那么 \(2^{\mathfrak{a}} \leqslant 1 + \mathfrak{a}\mathfrak{b}\) （注意，如果 \(\mathbf{Y} \subset \mathbf{A}\) 和 \(\mathbf{Y} \neq \emptyset\) ，那么 \(f(\mathbf{Y}) \in \mathbf{A}\) ).
\item
  设 B 是所有 \(x \in E\) 所成的集合，使得对每个非空子集 \(X\) （\(f^{-1}(x)\) 的子集）, 都有 Card \((X) \leqslant \mathfrak{b}\) . 证明 Card \((B) \leqslant \mathfrak{b}\) .
\end{enumerate}

\begin{enumerate}
\def\labelenumi{\arabic{enumi}.}
\setcounter{enumi}{4}
\tightlist
\item
  让 \((\lambda_{t})_{t\in I}\) 是序型族（§2，习题 13），由有序集 I 索引。证明
\end{enumerate}

\[
\operatorname{Card} \left(\sum_ {i \in I} \lambda_ {i}\right) = \sum_ {i \in I} \operatorname{Card} \left(\lambda_ {i}\right)
\]

并且，如果 I 是良序的，

\[
\operatorname{Card} \left(\underset {i \in I} {\mathbf {P}} \lambda_ {i}\right) = \underset {i \in I} {\mathbf {P}} \operatorname{Card} \left(\lambda_ {i}\right).
\]

\begin{enumerate}
\def\labelenumi{\arabic{enumi}.}
\setcounter{enumi}{5}
\tightlist
\item
  证明对于每个集合 E 都存在 \(\mathbf{X} \subset \mathbf{E}\) 使得 \(\mathbf{X} \notin \mathbf{E}\) （使用第 6 号的定理 2）。
\end{enumerate}

